\documentclass[11pt]{article}
\usepackage{amsmath,amssymb,amsthm,hyperref,geometry}\geometry{margin=1in}
\newtheorem{theorem}{Theorem}\newtheorem{proposition}{Proposition}
\title{Continuous completion of a derived display interaction}
\date{15 November 2026}\begin{document}\maketitle
\paragraph{Status.} Proposed general extension for independent review, conditional on the independently reviewed display-exact criterion. The finite interaction is derived first; the completion conditions below do not assume continuity of its mixed crossing. Final-sequence arguments are established categorical machinery, not a novelty claim. This note does not establish a total trace or an arbitrary infinite probability law.

\section{Continuity on the independent presentation}
Fix the independent display context and content data of the display-exact criterion. Write $\operatorname{Arr}_{\mathcal D}(\mathcal C)$ for the full category of admitted display arrows. Impose the following additional, separately checkable conditions for inverse $\omega$-chains.
\begin{enumerate}
\item $\mathcal C$ and the total content category $L$ have such limits; $p,d$ preserve them. The display-arrow subcategory is closed under limits calculated in $\operatorname{Arr}(\mathcal C)$.
\item $S$ preserves such limits. No assertion that a finitary or Cartesian functor automatically does so is made.
\item The jointly indexed Cartesian-lift functor $(f,Q)\mapsto f^*Q$ and opcartesian-transport functor $(g,Q)\mapsto g_!Q$ preserve these limits on their actual typed domains. The domains are pullbacks of the arrow and grade categories, not a collection of unrelated fixed fibres. The chosen cleavages and their coherent comparisons may be replaced by equivalent split presentations if required; continuity means the canonical limit comparisons are isomorphisms.
\item The independent content unit and typed accumulation preserve these limits. Their typed domain is $\{(Q,R):pR=dQ\}$. Context unit and display composition have this property from the first condition.
\end{enumerate}
All categories of typed tuples below use these actual pullback domains (or their equivalent coherent pseudo-pullback presentations when preservation is specified up to isomorphism). Their limits are created componentwise by the hypotheses on the typing functors. The assumptions are deliberately stronger than existence of a limit of underlying spaces: a lost native grade, law or observation is not repaired by a limit of carriers.

\begin{theorem}[Continuity of the derived crossing]
On compatible chains of native represented inputs $(Q,\pi:E\to pQ,f:B\to dQ,g:E\to C)$, the full derived crossing commutes with inverse limits. Its four comparison isomorphisms and ten equations commute with the same limit comparisons. The assertion is on retained presented categories and hence applies to any compatible native representative chain; it does not assert that arbitrary pseudocoends commute with limits.
\end{theorem}
\begin{proof}
For a chain indexed by $n$, limits commute with pullbacks, so
\[
 \lim_n(dQ_n\times_{S pQ_n}SE_n)
 \cong d(\lim_nQ_n)\times_{S(\lim_npQ_n)}S(\lim_nE_n).
\]
This is precisely the carrier of the crossing of the limit, with its two projections. Closure of displays ensures it remains an admitted context. The output content is $g_!\pi^*Q$; its limit comparison is an isomorphism by joint continuity of the two independent transport functors. The body square is the unique map with the same display projections, followed by the natural opcartesian body, so it also commutes with the comparison. This proves the result on objects and native arrows, not merely extensions. Pure balance is respected because its canonical mate is natural; no exchange of coend and limit is used.

The four comparison maps were defined by pullback pasting, Cartesian monad unit/multiplication and Cartesian-grade factorizations. Each commutes with these limit comparisons by its universal characterization and the independent continuity assumptions. Every one of the ten equations therefore remains an equality after passage to the limit. More directly, project the proposed equation to every chain stage; it is the already proved finite equation there. Limit projections are jointly monic on arrows. For grades this argument is in $L$, not after the possibly nonfaithful functor $d$.
\end{proof}

\section{Guarded constructions and actual completion}
The following standard theorem explains what the extra assumptions buy. It is stated at a general presentation category $\mathcal E$, which may be a typed category of contexts, contents or native bodies when its specified primitives act on that category. One must establish that typing and the existence of the terminal object in the chosen category; it is not inferred for every comma fibre.

\begin{theorem}[Final-sequence completion and native comparison]
Let $\mathcal E$ have a terminal object and inverse $\omega$-limits. Let $F:\mathcal E\to\mathcal E$ preserve the limit of its final sequence. Then $Z_F=\lim_nF^n1$ has a canonical final $F$-coalgebra structure. If $G$ satisfies the same conditions and $\kappa:F\Rightarrow G$ is a native natural transformation, the finite comparison maps extend uniquely to a map $Z_F\to Z_G$. This assignment preserves identities and vertical composition of natural transformations.
\end{theorem}
\begin{proof}
Let $b_n:F^{n+1}1\to F^n1$ be the final-sequence bonds. Preservation gives $FZ_F\cong\lim_nF^{n+1}1\cong Z_F$, with the latter comparison obtained by deleting the terminal first stage. Invert to obtain $z_F:Z_F\to FZ_F$. For a coalgebra $c:X\to FX$, define $c_0:X\to1$ and $c_{n+1}=Fc_n\,c$. These form a cone; its unique map into $Z_F$ is a coalgebra morphism because equality can be tested at every projection. Conversely every coalgebra morphism has these projections, proving finality.

Put $u_0=\mathrm{id}_1$ and $u_{n+1}=G u_n\,\kappa_{F^n1}$. Naturality shows these commute with the bonds. Their limit defines the comparison; equivalently it is coiteration of $\kappa_{Z_F}z_F$. Identity and composition follow either by induction on $n$ using naturality or by joint monicity of limit projections. None of these maps is claimed invertible unless additional hypotheses imply it.
\end{proof}

For example-free use, specify a guarded one-step endofunctor $F$ as a finite expression in the admitted typed primitive functors, finite limits and constant parameters. If the expression is genuinely an endofunctor on $\mathcal E$ and each primitive is continuous there, then $F$ is continuous: finite limits commute with inverse limits, and compositions of continuous functors are continuous. Derived crossing may be used in such an expression by the preceding theorem, not as an unproved continuous primitive. Natural native comparisons generated by the interaction law then extend to these completed recursive objects. Equality of two such generated comparisons follows from equality at all finite stages. This supplies recursive transport for the specified guarded constructions; it is not a classification of every recursive syntax or every feedback equation.

\section{An actual recurrent context construction}
This section specifies a recursive endofunctor directly from the interaction doctrine, rather than leaving the preceding theorem without a typed application. Suppose additionally that $\mathcal C$ has a terminal object. Fix $Q\in L_A$ and an explicit recurrent-boundary display $r:dQ\to A$. This return-to-boundary arrow is extra input; it is not supplied automatically by an arbitrary content grade. Define on the actual context fibre $K_A=\mathcal D/A$
\[
 F_{Q,r}(\pi:E\to A)
   =\bigl(dQ\times_{SA}SE\xrightarrow{\mathrm{pr}_1}dQ\xrightarrow r A\bigr).
\tag{R}
\]
On native context arrows use the pullback map induced by the identity of $dQ$ and $S$ of that arrow. Thus (R) is precisely the constructed crossing context followed by display accumulation, not a new supplied feedback operator.

\begin{theorem}[Recurrent display unfolding]
Under the stated continuity conditions, $F_{Q,r}$ has a final coalgebra in $K_A$. Its carrier is an admitted display over $A$, obtained as the inverse limit of its finite unfolding contexts. A native content arrow $b:Q\to Q'$ over $\mathrm{id}_A$ satisfying $r'db=r$ induces a natural transformation $F_{Q,r}\Rightarrow F_{Q',r'}$ and hence a directed map between their final unfolding contexts. These maps preserve identities and composition.
\end{theorem}
\begin{proof}
The fibre $K_A$ has terminal object $\mathrm{id}_A$. Its inverse-chain limits are calculated in $\mathcal C/A$ and remain displays by closure; nonempty connected inverse limits commute with the constant base $A$. The maps in (R) are displays because $S\pi$ and $r$ are displays. Its formation preserves inverse limits by continuity of $S$, commutation of limits with pullbacks, and the explicit slice construction. Thus the final-sequence theorem applies to this actual category and functor. Its structure map retains the $dQ$ witness projection and the $SE$ continuation projection together; neither is discarded in the limit.

Naturality of $\nu$ gives $\nu_{Q'}db=\nu_Q$. Therefore $(db,\mathrm{id}_{SE})$ induces a map between the two pullbacks in (R), and $r'db=r$ makes it an arrow of $K_A$. Pullback uniqueness proves naturality and compatibility with identity/composition. The final-sequence comparison theorem supplies the asserted infinite native maps.
\end{proof}

This is a categorical dependent-unfolding theorem under explicit continuity. It does not turn the actual return arrow $r$ into an arbitrary feedback wire, prove a trace axiom, or invent an infinite effect from finite marginals. If the context category carries additional semantic data, the limits must be limits there with the stated compatibility, not merely limits of forgotten carriers. No set-theoretic enumeration of occurrences or probabilistic instantiation is a premise of the construction.

\section{Without completion, and why delay is not silently interchanged}
Even without limits in $\mathcal E$, the final sequence and its native comparison exist as actual inverse diagrams whenever a terminal object exists. They provide finite-stage semantics, not an infinite object in $\mathcal E$. They do not provide an infinite effect by taking a limit in some different category of bare sets.

There is a separate typing issue with syntactic delay. On inverse diagrams the usual delay sets $(\triangleright X)_0=1$ and $(\triangleright X)_{n+1}=X_n$. A pointwise functor $S$ need not commute with delay: at stage zero the two sides are $S1$ and $1$. In particular Cartesianness of an occurrence monad does not imply $S1\cong1$. The completion theorem does not assume or use such a commuting isomorphism. A full delay-distributive interaction or guarded trace requires its own appropriately typed structure and laws. This is a limitation of that proposed interchange, not an impossibility of feedback in general.

\paragraph{Background.} The final-sequence construction is standard coalgebra theory; see J. Ad\'amek, \emph{Introduction to coalgebra}, and the explicit final-sequence theorem recalled in \emph{Coinductive Streams in Monoidal Categories}, \url{https://lmcs.episciences.org/16337/pdf}. Guarded inverse-diagram models have a separate established literature, e.g. Birkedal, M\o gelberg, Schwinghammer and St\o vring, \emph{First steps in synthetic guarded domain theory}, \url{https://lmcs.episciences.org/1041}. The project-specific obligation here is continuity and native comparison of the independently constructed interaction, not a claim to originate these completion methods.
\end{document}
